% Unit 35 — Modified famous integrals
\input{technique-note.tex}
\title{Integration Bee --- Unit 35:\\ Modified Famous Integrals}
\author{UK Uni Integration Bee prep\\ Scale a named integral, or expand a series}
\date{}
\begin{document}
\maketitle
\section*{Status}
Tags: \texttt{yang:6.8}; techniques: \texttt{modified-famous}.
A famous integral with one constant changed is still that integral after a substitution. Yang \S6.8 collects those tweaks. The values come from unit 30, from the geometric series, and from one double-angle argument for $\ln(\sin x)$. None of these drills needs a new constant invented on the spot.
\section{Official anchors}
\begin{itemize}
\item Unit 30 for the Gaussian and $\int_0^{\infty}(\sin x)/x\,dx=\pi/2$.
\item The Basel value $\sum_{n=1}^{\infty}1/n^2=\pi^2/6$, and the alternating partner $\sum_{n=1}^{\infty}(-1)^{n-1}/n^2=\pi^2/12$.
\item Yang Topics Guide 2nd Ed \S6.8 Modified Famous Integrals. Syllabus \S2 names zeta.
\end{itemize}
\section{Learning goals}
\begin{enumerate}
  \item Scale the limit and the constant in the exponent or the argument. Write the substitution.
  \item Expand $1/(1\pm x)$ and integrate $x^n\ln x$ termwise when the famous value is a zeta number.
\end{enumerate}

\section{Scale the Gaussian --- \texttt{U35-D01}}
\[
\int_0^{\infty}e^{-2x^2}\,dx=\frac12\sqrt{\frac{\pi}{2}}=\frac{\sqrt{2\pi}}{4}.
\]
Put $u=x\sqrt{2}$ in \texttt{U30-D01}. \texttt{U35-D03} is $\int_0^{\infty}e^{-3x}\,dx=1/3$, which is elementary and is not a Gaussian.

\section{Scale the Dirichlet integral --- \texttt{U35-D02}}
\[
\int_0^{\infty}\frac{\sin 2x}{x}\,dx=\frac\pi2.
\]
Put $u=2x$. Then $dx/x=du/u$, and the integral is \texttt{U30-D05}. The factor $2$ in the sine does not change the value.

\section{Series --- \texttt{U35-D04} and \texttt{U35-D05}}
On $[0,1)$,
\[
\frac{1}{1-x}=\sum_{n=0}^{\infty}x^n,\qquad \int_0^1 x^n\ln x\,dx=-\frac{1}{(n+1)^2}.
\]
So
\[
\int_0^1\frac{\ln x}{1-x}\,dx=-\sum_{m=1}^{\infty}\frac{1}{m^2}=-\frac{\pi^2}{6}.
\]
The alternating expansion of $1/(1+x)$ produces
\[
\int_0^1\frac{\ln x}{1+x}\,dx=-\sum_{m=1}^{\infty}\frac{(-1)^{m-1}}{m^2}=-\frac{\pi^2}{12}.
\]
Those are the stored answers. The endpoint $x=1$ in \texttt{U35-D04} is improper. The series of positive terms and the integral agree as improper integrals. \texttt{U35-D08} is the elementary cousin $\int_0^1 -\ln x\,dx=1$.

\section{Logarithm of sine --- \texttt{U35-D06}}
\[
I=\int_0^{\pi/2}\ln(\sin x)\,dx.
\]
The reflection $x\mapsto\pi/2-x$ shows $I=\int_0^{\pi/2}\ln(\cos x)\,dx$, so
\[
2I=\int_0^{\pi/2}\ln(\sin x\cos x)\,dx=\int_0^{\pi/2}\ln(\sin 2x)\,dx-\int_0^{\pi/2}\ln 2\,dx.
\]
The second integral is $(\pi/2)\ln 2$. For the first, $v=2x$ gives $\frac12\int_0^{\pi}\ln(\sin v)\,dv=\int_0^{\pi/2}\ln(\sin v)\,dv=I$, because $\sin(\pi-v)=\sin v$. Therefore $2I=I-(\pi/2)\ln 2$, and $I=-(\pi/2)\ln 2$. \texttt{U36-D02} is the same integral.

\texttt{U35-D07} is $(e^{-x}-e^{-2x})/x$ on $[0,\infty)$, stored value $\ln 2$. That is \texttt{U34-D01}. Use the parameter argument there. Do not start a series.

\section{Drill pointers}
See \texttt{drills.tex} / \texttt{drills.md}. Feed: \texttt{verify/drills\_unit35.json}.

\section*{On a miss}
Teach this file. Name the famous integral you scaled, or the series you integrated. Do not invent an answer that is not in \texttt{drills.md}.
\end{document}
